%======================================================================
\section{Formal verification and its scope}\label{app:formal}
%======================================================================

The Lean sources are in the \texttt{lean/} directory of the accompanying
repository. The exact types, dependency pins, source hashes, third-party notices
and commands for a separate build and kernel replay are given in
\texttt{lean/README.md}, the dependency record and the source manifest
in that directory. The frozen file \texttt{Challenge.lean}
contains historical admitted declarations; these are statement markers and are
not the proofs of the certified roots. The final \texttt{Audit.lean} compares
its six analytic statement types with the proved declarations and prints the
axioms of every listed root. The A/E/F aggregate propositions are defined in
\texttt{Blueprint.lean}, beyond the original frozen statement file.

{\small
\setlength{\tabcolsep}{4pt}
\begin{longtable}{@{}p{0.34\textwidth}p{0.61\textwidth}@{}}
\toprule
Paper statement & Lean declaration (qualified by namespace)\\
\midrule
\endhead
Theorem A & \texttt{Blueprint.theoremA}\\
\Cref{lem:conjugator-length} & \texttt{Blueprint.DiagramProof.}\newline\texttt{conjugatorLengths}\\
A second Theorem A wrapper & \texttt{Blueprint.DiagramProof.theoremA}\\
Pauli hash models used in the upper bound & \texttt{Blueprint.Published.pauliHashModels}\\
\Cref{thm:count} & \texttt{Companion.Analytic.CountProof.count}\\
Theorem C and \cref{cor:count} & \texttt{Companion.Analytic.CountProof.}\newline\texttt{uniformCount}\\
\Cref{thm:transfer} & \texttt{Companion.Analytic.TransferProof.}\newline\texttt{transfer}\\
\Cref{thm:return}, including the specified measurement and zero error & \texttt{Companion.Analytic.ReturnProof.}\newline\texttt{returnTheorem}\\
Positive-error existence from \cref{thm:return} & \texttt{Companion.Analytic.ReturnProof.}\newline\texttt{returnExistence}\\
Theorem D and the three conclusions of \cref{cor:two-sided} & \texttt{Companion.Analytic.ProfileProof.}\newline\texttt{twoSided}\\
Theorem E & \texttt{Blueprint.theoremE}\\
Theorem F & \texttt{Blueprint.theoremF}\\
\Cref{thm:sector}, entropy and rank inequalities & \texttt{Blueprint.ASector.source\_sector}\\
\Cref{prop:hash-models} & \texttt{Blueprint.AHash.general\_hash\_models}\\
Literal boundary merge in the conjugator proof & \texttt{Blueprint.DiagramProof.merge\_matching}\\
\bottomrule
\end{longtable}}

\paragraph{Exact scope.}
The type of \texttt{Blueprint.theoremA} is
\texttt{Published.TensorPerfectFinite} $\to$ \texttt{AStatement}, in namespace
\texttt{Blueprint}. The latter includes the concrete counts, arity three,
quantum value one, exclusion of finite perfect strategies, pure and mixed lower
bounds, the maximally entangled upper bound, and dimension-profile exponents.
It additionally includes exclusion of perfect spatial strategies; that conclusion
uses the published parameter described above. The pure lower bound also bounds
entropy by log rank. The entropy asymptotic is supplied by the lower and upper
bounds rather than a separate entropy-limit conjunct.
The uniform count includes the rank conclusion. The analytic statements are
finite-dimensional; linear-system rows are nonempty, sorted, duplicate-free and
bounded, with an admissible question distribution. The return theorem fixes
Bob's rounded observables and Alice's sequential measurement with parity
correction. The sector and hash-model adapter declarations retain, respectively,
a conjugator-length and a Pauli-hash-model parameter; the Theorem~A assembly
supplies both proofs internally. The sector root records the entropy and rank
inequalities, not all intermediate sector-state conclusions as separate conjuncts.
The literal boundary merge is an internal combinatorial lemma, not an additional
paper theorem. The E/F propositions use the definitions of profiles in
\texttt{Blueprint/RatesDefinitions.lean} and \texttt{Blueprint/FSVGame.lean}.
The K profile has the explicit range $0<\delta\le2$.

\paragraph{What is not formalized.}
The external perfect-tensor-strategy implication itself is a parameter rather
than a Lean proof. There is no standalone certified root for Theorem~B, no claim
that every proposition, discussion, literature result or appendix of the paper
has been formalized, and no formal certification of the ancillary numerical
experiments, the prose, or its identification with the definitions. Historical
comments in the unchanged statement and proof sources may describe an earlier
state of completion; the proved declarations and their audited types specify
the present claims. In particular, the remaining point of trust is that the
definitions in \texttt{Challenge.lean} mean what the paper means.
